\documentclass[12pt]{article}

\usepackage[a4paper, total={6in, 8in}]{geometry}
\usepackage{textcomp}

\begin{document}
\setlength{\parindent}{0pt}

\def \t {\textrightarrow}
\def \v {\vspace{1em}}
\def \bi {\begin{itemize}}
\def \ei {\end{itemize}}
\def \s[#1] {\section*{#1}}
\def \ss[#1] {\subsection*{#1}}
\def \sss[#1] {\subsubsection*{#1}}

\s[The 1920s - 1930s]
In the USA the 1920s were a decade of prosperity (in UK they were hard times) \t these years are known as "the roaring twenties" \\
A novel which reflects on the execess of this age is "The great Gatsby" by Fitzgerald \t the main charcters are young rich people, but they are unable to reach happiness because they have to go through personal tragedies \\
These years ended in 1929 with the Wall Street crisis \t so 1930s where hard times, until 1935 with the New Deal

\v

But age of anxiety \t ones mean the intellectual atmosphere that characterizes these years in the USA \t most american intellectuals felt axienty, due to the loss of certainties \\
That collective loss were caused by:
\bi
  \item the great war \t it was a collective traumatic experience:
    \bi
      \item 10 million people died, mainly soldiers but also many civilians
      \item it showed how cruel humans can be
      \item shows how progress can be detrimental and devastating, if not controlled by ethics
    \ei
  \item revolutionary theories \t either elaborated in those years, or elaborated before but became accepted in the 20s and 30s \t they contradicted many ideals those intellectuals believed in:
    \bi
      \item Darwin's theory \t after WWI became actually accepted \t it denies years of biblical education, and humans are not privileged creatures
      \item Freud's theory \t in particual his thre part model of the human psyche \t most of the psyche is dark and unexplored:
        \bi
          \item id: unconscious, impulsive
          \item ego: the rational side
          \item super ego: moral norms humans are taught to live in society
        \ei
      \item Einstein's theory \t theory of relativity, which made the concepts of time and space relative and not absolute
      \item Bergson's theory \t he elaborates a theory of time, according to which there are two types of time:
        \bi
          \item external objective time, measured by clocks
          \item psycholigical, subjective time \t called duration
        \ei
      \item Frazer's theory \t he is an anthropologist (science that originated in the second part of the 19th century, and studies human societies all around the world) \t in particular primitive societies \t in "The Golden Bough" he writes about his studies regarding some african societies, and says that even if they are primitive, these societies are complex and articulated = discovery clashes with imperialism \t now some intellectuals (modernists) got interested in the culture of these populations
    \ei
\ei

\ss[The second coming - William Butler Yeats]
He is a modernist irish poet \\
By reading this poem, we can understand what it meant to be an intellectual during the age of anxiety \\
This poem was written in 1919 \t the war had just finished and the europeans economies were devastated, + there the russian revolution just happened \\
Between 1919 and 1921 there was a civil war in Ireland, which concluded with the creation of the Free Republic of Ireland \t then another one started in northern ireland ?? \\
Yeats elaborated a personal theory of history \t he believed that human history was divided into cycles lasting 2000 years each \\
He also gave a graphic representation of this cycles \t each one has the shape of gyre

\v

Two stanzas \t the first one is objective (it presents a description of the world at the time) and shorter \t the second subjective because it is about a vision of the future and longer \\
He was leaving in 1919, ending of the cycle that started with Christ's birth \t the falcon cannot hear the falconer stands for this and is a metaphore \t this means that the cycle is about to finish \\
Turning \t is a diacope \\
Falcon and falconer is a polyptoton \\
Mere anarchy \t there is not a political orientation in that world \\
The tide \t refernce of the great world war \\
The world functions the other way around

\v

There is a vision about the future, that lyrical eye has \\
The first coming was the birth of christ, now the new cycle begins with a second coming \t it has not happened yet, he has a vision about the future \\
As the sun in the desert \t simile \\
The new creature is monstrous one, a sphinx-like creature \\
Rocking cradle = culla ondeggiante \\
Jesus was born in Bethlem \t the creature is anit-christ

\v

The message is that the historical cycle is finishing, and a new one is going to start \t and nightmarish one \\
Yeats was an irish independentits \t he also wrote political poems

\s[Anglo modernism]
Definition: \textit{A cultural movement developing in the ENglish-speaking contries starting to emerge at the end of the 19th century, more clearly after the Great War, as a reaction against the wrongs of its age with its puritanical, authoritative, imperialistic and materialistic attitude. Modernism likwise reacted against the art of the past, considere as unable to represent and understand the modern world} \\
High modernism is the one that was born after the WWI \\
Modernist artists tried to find new ways to understand reality \\
The one before WWI went against past art \t after they recognize the value of the art of the past, but this art has to be adapted
\end{document}
